% ==========================================
% Archivo: sections/5_conclusions.tex
% ==========================================
\section{Conclusiones}

Del análisis de las trayectorias balísticas estudiadas se concluye que el efecto del rozamiento domina frente a las variaciones del campo gravitatorio. Cuantitativamente, el rozamiento induce desviaciones respecto a la trayectoria ideal del orden de $10^3 - 10^4 \, \mathrm{m}$, mientras que la corrección por gravedad dependiente de la altura produce perturbaciones de segundo orden ($\Delta \sim 10^0 - 10^2 \, \mathrm{m}$). Esta dominancia disipativa altera el ángulo de lanzamiento que maximiza el alcance, reduciendo el ángulo óptimo desde los $\theta = 45^\circ$ del caso ideal hasta $\theta = 38.8^\circ$ en presencia de aire. En cuanto a la introudcción de los modelos isotérmicos y adiabáticos, estos modifican dicho ángulo óptimo a $\theta = 47.2^\circ$ y $\theta = 43.7^\circ$ para los casos isotérmico y adiabático, respectivamente.

Respecto al efecto Magnus, se demuestra que la topología de la trayectoria presenta una fuerte dependencia con la dirección del vector velocidad angular $\boldsymbol{\omega}$ debido al acoplamiento de las ecuaciones de movimiento. Cuando la rotación es perpendicular a la gravedad ($\boldsymbol{\omega} \perp \vec{g}$), se genera una fuerza de sustentación vertical que gobierna la dinámica, permitiendo alcanzar distancias extremas ($113\, \mathrm{km}$ para un ángulo de lanzamiento de $\theta = 6.4^\circ$) y propiciando la aparición de \textit{loops} en la trayectoria. Por el contrario, si el eje de rotación es paralelo a la gravedad ($\boldsymbol{\omega} \parallel \vec{g}$), fenómeno característico de los efectos \textit{slice} y \textit{hook}, la perturbación actúa de forma puramente transversal. Esto rompe la simetría en el plano $xy$, induciendo desviaciones laterales ($\Delta y = \pm 46.7 \, \mathrm{m}$) estrictamente simétricas bajo la inversión de $\boldsymbol{\omega}$, pero preservando la componente parabólica vertical al anularse el acoplamiento con $\vec{g}$.

Por último, el análisis de los sistemas multicuerpo acoplados mediante muelles viscoelásticos ilustra de forma clara la separación entre el movimiento del centro de masas y la dinámica interna. Al no existir fuerzas disipativas externas, el centro de masas sigue estrictamente una trayectoria parabólica balística. La fricción viscoelástica de los muelles provoca una rápida transición del régimen transitorio inicial a soluciones estacionarias: un patrón helicoidal en el caso de dos cuerpos, y una trayectoria cuasiperiódica en el caso de tres cuerpos.
